\documentclass[12pt,openright,oneside,a4paper]{report}

% -------------------------------------------------
% PACOTES E DIRETRIZES DA ABNT
% -------------------------------------------------
\usepackage[utf8]{inputenc}
\usepackage[T1]{fontenc}
\usepackage[brazil]{babel}
\usepackage{indentfirst}
\usepackage{microtype}
\usepackage{setspace}
\usepackage{graphicx}

% Margens Oficiais da ABNT
\usepackage[top=3cm, bottom=2cm, left=3cm, right=2cm]{geometry}

% Fonte Arial / Helvetica para sumário e títulos limpos
\usepackage{helvet}
\renewcommand{\familydefault}{\sfdefault}

\onehalfspacing
\setlength{\parindent}{1.25cm}
\setlength{\parskip}{0cm}

% -------------------------------------------------
% DADOS DA PUBLICAÇÃO
% -------------------------------------------------
\title{Revista 2026/1}
\author{Universidade Federal de São Paulo - UNIFESP}
\date{2026}

\begin{document}

% -------------------------------------------------
% CAPA
% -------------------------------------------------
\begin{titlepage}
    \centering
    \vspace*{2cm}
    {\Large UNIVERSIDADE FEDERAL DE SÃO PAULO - UNIFESP \par}
    \vspace{0.5cm}
    {\large PROGRAMA PEPICT \par}
    \vfill
    {\bfseries\LARGE Revista 2026/1 \par}
    \vspace{1cm}
    {\large Revista Institucional de Relatórios Acadêmicos \par}
    \vfill
    {\large São José dos Campos\par}
    {\large 2026 \par}
\end{titlepage}

% -------------------------------------------------
% FOLHA DE ROSTO
% -------------------------------------------------
\begin{titlepage}
    \centering
    {\Large UNIVERSIDADE FEDERAL DE SÃO PAULO - UNIFESP \par}
    \vspace{0.5cm}
    {\large PROGRAMA PEPICT \par}
    \vspace{4cm}
    {\bfseries\LARGE Revista 2026/1 \par}
    \vspace{2cm}
    \hspace{.45\textwidth}
    \begin{minipage}{.5\textwidth}
        \onehalfspacing
        Revista institucional destinada à publicação de relatórios acadêmicos desenvolvidos por alunos vinculados às unidades curriculares e projetos do programa PEPICT.
    \end{minipage}
    \vfill
    \begin{center}
        São José dos Campos \\
        2026
    \end{center}
\end{titlepage}

% -------------------------------------------------
% EXPEDIENTE INSTITUCIONAL
% -------------------------------------------------
\chapter*{Informações da publicação}
\thispagestyle{empty}

\begin{flushleft}
{\large \textbf{UNIVERSIDADE FEDERAL DE SÃO PAULO – UNIFESP}}
\vspace{1cm}

\textbf{Programa PEPICT} \\
NOME_COORDENADOR_GERAL_PEPICT
\vspace{0.7cm}

\textbf{Coordenação Editorial} \\
NOME_COORDENADOR_EDITORIAL
\vspace{0.7cm}

\textbf{Comissão Organizadora} \\
NOME_COMISSAO
\vspace{0.7cm}

\textbf{Revisão} \\
NOME_REVISAO
\vspace{0.7cm}

\textbf{Diagramação} \\
NOME_DIAGRAMACAO
\vspace{1.5cm}

\textbf{\large Revista 2026/1} \\
2026
\vspace{0.5cm}

ORGANIZADORES
\end{flushleft}
\clearpage

% -------------------------------------------------
% SUMÁRIO
% -------------------------------------------------
\tableofcontents
\clearpage

% -------------------------------------------------
% INTRODUÇÃO
% -------------------------------------------------
\chapter{Introdução}
_INTRODUCAO_REVISTA_
\clearpage

% -------------------------------------------------
% CONTEÚDO DINÂMICO DOS CAPÍTULOS
% -------------------------------------------------
\chapter{Arquitetura e Organização de Computadores}

\section{Projeto - Panda}
\noindent
\textbf{Autores}
Ana Carolina

\vspace{0.7cm}
\noindent
\textbf{Disciplina e Curso}
Engenharia de Computação

\vspace{0.7cm}
\noindent
\textbf{Programa PEPICT vinculado}
Tecnologia e Inovação para o Desenvolvimento Social

\vspace{0.7cm}
\noindent
\textbf{Objetivos}
muitos

\vspace{0.7cm}
\noindent
\textbf{Metodologia}
o que eu quiser

\vspace{0.7cm}
\noindent
\textbf{Resultados e Produtos}
funcionou

\vspace{0.7cm}
\noindent
\textbf{Relação com os ODS}
é aquilo né

\vspace{0.7cm}
\noindent
\textbf{Reflexão Crítica}
refletido

\vspace{0.7cm}
\noindent
\textbf{Referências}
ual

\vspace{0.7cm}
\section{Projeto - Panda}
\noindent
\textbf{Autores}
Aluno teste2

\vspace{0.7cm}
\noindent
\textbf{Disciplina e Curso}
Engenharia de Computação

\vspace{0.7cm}
\noindent
\textbf{Programa PEPICT vinculado}
Educação e Sustentabilidade

\vspace{0.7cm}
\noindent
\textbf{Objetivos}


\vspace{0.7cm}
\noindent
\textbf{Metodologia}


\vspace{0.7cm}
\noindent
\textbf{Resultados e Produtos}


\vspace{0.7cm}
\noindent
\textbf{Relação com os ODS}


\vspace{0.7cm}
\noindent
\textbf{Reflexão Crítica}


\vspace{0.7cm}
\noindent
\textbf{Referências}
None

\vspace{0.7cm}
\section{titulo do trabalho}
\noindent
\textbf{Autores}
Aluno Teste2

\vspace{0.7cm}
\noindent
\textbf{Disciplina e Curso}
Engenharia de Computação

\vspace{0.7cm}
\noindent
\textbf{Programa PEPICT vinculado}
Educação e Sustentabilidade, Tecnologia e Inovação para o Desenvolvimento Social

\vspace{0.7cm}
\noindent
\textbf{Objetivos}
sistema deve permitir que o coordenador filtre turmas por semestre, visualize projetos aprovados, selecione os trabalhos que farão parte da revista e gere o arquivo PDF consolidado da edição.

\vspace{0.7cm}
\noindent
\textbf{Metodologia}
Concluída a etapa de requisitos, o passo seguinte foi estruturar o modelo de resumo. Para isso, realizamos reuniões com a orientação do projeto para alinhar e selecionar quais informações seriam essenciais e pertinentes para compor esses resumos estruturados, o que resultou no modelo ilustrado na imagem abaixo:

\vspace{0.7cm}
\noindent
\textbf{Resultados e Produtos}
Foi desenvolvida em 2001, quando seus autores buscaram solucionar lacunas e ineficiências do desenvolvimento em cascata, detectadas quando essa metodologia não produzia mais resultados satisfatórios frente a vários parâmetros de medida. Era comum, por exemplo,  demorar anos entre a identificação de uma necessidade da empresa e a entrega do sistema pronto. Nesse meio-tempo, as mudanças no mercado e nas exigências do negócio faziam com que grande parte do projeto fosse cancelada antes mesmo de ser entregue. Esse desperdício de tempo e dinheiro motivou vários desenvolvedores a buscarem uma alternativa.

\vspace{0.7cm}
\noindent
\textbf{Relação com os ODS}
A metodologia ágil é uma abordagem de desenvolvimento que visa criar um fluxo de trabalho para tornar o processo de produção o mais eficiente e assertivo possível. Nessa abordagem, o desenvolvimento é um processo colaborativo no qual todas as partes mantêm um diálogo claro, a fim de manter o produto o mais próximo do desejado pelo solicitante, dentro da realidade de desenvolvimento e recursos disponíveis

\vspace{0.7cm}
\noindent
\textbf{Reflexão Crítica}
A metodologia ágil é uma abordagem de desenvolvimento que visa criar um fluxo de trabalho para tornar o processo de produção o mais eficiente e assertivo possível. Nessa abordagem, o desenvolvimento é um processo colaborativo no qual todas as partes mantêm um diálogo claro, a fim de manter o produto o mais próximo do desejado pelo solicitante, dentro da realidade de desenvolvimento e recursos disponíveis

\vspace{0.7cm}
\noindent
\textbf{Referências}
A metodologia ágil é uma abordagem de desenvolvimento que visa criar um fluxo de trabalho para tornar o processo de produção o mais eficiente e assertivo possível. Nessa abordagem, o desenvolvimento é um processo colaborativo no qual todas as partes mantêm um diálogo claro, a fim de manter o produto o mais próximo do desejado pelo solicitante, dentro da realidade de desenvolvimento e recursos disponíveis

\vspace{0.7cm}
\chapter{Sistemas Embarcados}

\section{reenvio}
\noindent
\textbf{Autores}
Ana Carolina

\vspace{0.7cm}
\noindent
\textbf{Disciplina e Curso}
Engenharia de Computação

\vspace{0.7cm}
\noindent
\textbf{Programa PEPICT vinculado}
Educação e Sustentabilidade

\vspace{0.7cm}
\noindent
\textbf{Objetivos}
Objetivos do projeto

\vspace{0.7cm}
\noindent
\textbf{Metodologia}
metodologia do projeto

\vspace{0.7cm}
\noindent
\textbf{Resultados e Produtos}
resultados do projeto

\vspace{0.7cm}
\noindent
\textbf{Relação com os ODS}
porque sim

\vspace{0.7cm}
\noindent
\textbf{Reflexão Crítica}
A definir

\vspace{0.7cm}
\noindent
\textbf{Referências}
aquelas

\vspace{0.7cm}
\section{Novo Projeto}
\noindent
\textbf{Autores}
eu

\vspace{0.7cm}
\noindent
\textbf{Disciplina e Curso}
Engenharia de Computação

\vspace{0.7cm}
\noindent
\textbf{Programa PEPICT vinculado}
Educação e Sustentabilidade, Tecnologia e Inovação para o Desenvolvimento Social

\vspace{0.7cm}
\noindent
\textbf{Objetivos}
obsetivos teste

\vspace{0.7cm}
\noindent
\textbf{Metodologia}
metodologia teste

\vspace{0.7cm}
\noindent
\textbf{Resultados e Produtos}
resultados teste

\vspace{0.7cm}
\noindent
\textbf{Relação com os ODS}
osd teste

\vspace{0.7cm}
\noindent
\textbf{Reflexão Crítica}
refleti bastante

\vspace{0.7cm}
\noindent
\textbf{Referências}
SILVA, Maurício Samy. Criando sites com HTML: sites de alta qualidade com HTML e CSS . Novatec Editora, 2008.

\vspace{0.7cm}
\section{Sistema de detecção de parede}
\noindent
\textbf{Autores}
carolina Silva

\vspace{0.7cm}
\noindent
\textbf{Disciplina e Curso}
Engenharia de Computação

\vspace{0.7cm}
\noindent
\textbf{Programa PEPICT vinculado}
Educação e Sustentabilidade, Tecnologia e Inovação para o Desenvolvimento Social

\vspace{0.7cm}
\noindent
\textbf{Objetivos}
Aqueles objetivos lá

\vspace{0.7cm}
\noindent
\textbf{Metodologia}
Aquela mesma

\vspace{0.7cm}
\noindent
\textbf{Resultados e Produtos}
Deu certo mesmo

\vspace{0.7cm}
\noindent
\textbf{Relação com os ODS}
Achei que sim

\vspace{0.7cm}
\noindent
\textbf{Reflexão Crítica}
Refleti bastante

\vspace{0.7cm}
\noindent
\textbf{Referências}
minha cabeça

\vspace{0.7cm}
\section{Projeto - Embarcados}
\noindent
\textbf{Autores}
Gertrudes Felesmina

\vspace{0.7cm}
\noindent
\textbf{Disciplina e Curso}
Engenharia de Computação

\vspace{0.7cm}
\noindent
\textbf{Programa PEPICT vinculado}
Educação e Sustentabilidade

\vspace{0.7cm}
\noindent
\textbf{Objetivos}


\vspace{0.7cm}
\noindent
\textbf{Metodologia}


\vspace{0.7cm}
\noindent
\textbf{Resultados e Produtos}


\vspace{0.7cm}
\noindent
\textbf{Relação com os ODS}


\vspace{0.7cm}
\noindent
\textbf{Reflexão Crítica}


\vspace{0.7cm}
\noindent
\textbf{Referências}
None

\vspace{0.7cm}
\section{Titulo do projeto}
\noindent
\textbf{Autores}
Aluno Teste2,Aluno 2

\vspace{0.7cm}
\noindent
\textbf{Disciplina e Curso}
Engenharia de Computação

\vspace{0.7cm}
\noindent
\textbf{Programa PEPICT vinculado}
Educação e Sustentabilidade, Tecnologia e Inovação para o Desenvolvimento Social

\vspace{0.7cm}
\noindent
\textbf{Objetivos}
objetivos teste

\vspace{0.7cm}
\noindent
\textbf{Metodologia}
metodologia teste

\vspace{0.7cm}
\noindent
\textbf{Resultados e Produtos}
resultados teste

\vspace{0.7cm}
\noindent
\textbf{Relação com os ODS}
justificativa teste

\vspace{0.7cm}
\noindent
\textbf{Reflexão Crítica}
reflexao teste

\vspace{0.7cm}
\noindent
\textbf{Referências}
referencia teste
referencia teste 2

\vspace{0.7cm}
\chapter{UC teste}

\chapter{UC teste 2}

\section{Titulo do projeto}
\noindent
\textbf{Autores}
Aluno teste2, aluno 2, aluno 3

\vspace{0.7cm}
\noindent
\textbf{Disciplina e Curso}
BCT

\vspace{0.7cm}
\noindent
\textbf{Programa PEPICT vinculado}
Educação e Sustentabilidade

\vspace{0.7cm}
\noindent
\textbf{Objetivos}
objetivo teste

\vspace{0.7cm}
\noindent
\textbf{Metodologia}
metodologia teste

\vspace{0.7cm}
\noindent
\textbf{Resultados e Produtos}
resultado teste

\vspace{0.7cm}
\noindent
\textbf{Relação com os ODS}
justifcativa teste

\vspace{0.7cm}
\noindent
\textbf{Reflexão Crítica}
reflexao teste

\vspace{0.7cm}
\noindent
\textbf{Referências}
referencia teste
referencia teste 2

\vspace{0.7cm}
\section{projeto PEPICT teste}
\noindent
\textbf{Autores}
Aluno Teste2, Fernando, Miguel

\vspace{0.7cm}
\noindent
\textbf{Disciplina e Curso}
BCT

\vspace{0.7cm}
\noindent
\textbf{Programa PEPICT vinculado}
Educação e Sustentabilidade, Tecnologia e Inovação para o Desenvolvimento Social

\vspace{0.7cm}
\noindent
\textbf{Objetivos}
O objetivo geral deste trabalho consiste no desenvolvimento e implantação de uma plataforma web
responsiva voltada à otimização das rotas de coleta de materiais recicláveis por catadores autônomos
na região periférica de São José dos Campos. A motivação do projeto parte da constatação da extrema
vulnerabilidade socioeconômica desses trabalhadores, que frequentemente percorrem longas distâncias
sem a garantia de encontrar resíduos triados. O sistema visa aproximar os moradores locais (geradores
de resíduos) e os catadores cadastrados, promovendo a valorização do trabalho ambiental urbano e
estimulando a cultura do descarte correto no ambiente universitário e em seu entorno.

\vspace{0.7cm}
\noindent
\textbf{Metodologia}
O trabalho foi estruturado em quatro etapas principais, pautadas pelos princípios da extensão
universitária dialógica:
Diagnóstico e Imersão: Realização de três reuniões presenciais com a Cooperativa de Catadores
local e visitas técnicas comunitárias para entender os gargalos da rotina de coleta e o nível de
letramento digital dos usuários.
Desenvolvimento Ágil: Construção do software utilizando a metodologia Scrum, focando em uma
interface simplificada e de alta acessibilidade (mobile-first), com validações semanais junto aos
representantes dos catadores.
Ação Extensionista e Interação com a Comunidade: Condução de oficinas de capacitação no
ICT/Unifesp para o uso da plataforma pelas famílias de catadores, além de uma campanha de
conscientização nos bairros vizinhos para que os moradores utilizassem o sistema para agendar
coletas.
Avaliação e Feedback: Período de testes assistidos em campo por duas semanas, monitorando a
usabilidade da ferramenta e coletando críticas para ajustes de layout e de fluxo.

\vspace{0.7cm}
\noindent
\textbf{Resultados e Produtos}
Como produto central, foi gerado o protótipo funcional da plataforma web denominado "ReciclaAção",
disponível de forma gratuita. Os resultados quantitativos e qualitativos alcançados incluem:
Protótipo de Software: Sistema web com mapa interativo de pontos de descarte gerados por
moradores e um painel de gerenciamento simplificado para os catadores.
Oficinas de Capacitação: Realização de 2 oficinas práticas no laboratório de informática do ICT,
certificando 14 catadores autônomos para o uso autônomo da tecnologia.
Relatório Técnico-Social: Documento detalhando os indicadores socioeconômicos levantados na
comunidade e o mapeamento de 45 novos pontos de coleta recorrente.
Impacto Inicial: Durante a fase piloto, registrou-se um aumento estimado de 25% na renda semanal
dos catadores participantes devido à diminuição do tempo gasto em trajetos improdutivos.

\vspace{0.7cm}
\noindent
\textbf{Relação com os ODS}
ODS 8 - Trabalho Decente e Crescimento Econômico: O projeto contribui diretamente ao fornecer
uma ferramenta tecnológica que eleva a produtividade física e a rentabilidade financeira dos catadores
autônomos, diminuindo a penosidade do trabalho de rua e inserindo-os ativamente na cadeia de valor
da economia circular regional.
ODS 11 - Cidades e Comunidades Sustentáveis: A iniciativa fomenta a gestão inclusiva e sustentável
dos resíduos sólidos urbanos de São José dos Campos, transformando a relação da comunidade
acadêmica e civil com o descarte e destinação correta do lixo reciclável.
ODS 12 - Consumo e Produção Responsáveis: Através das campanhas extensionistas e do uso da
plataforma, houve uma expressiva redução do volume de resíduos passíveis de reciclagem que seriam
erroneamente direcionados para

\vspace{0.7cm}
\noindent
\textbf{Reflexão Crítica}
A vivência na extensão universitária proporcionada por este projeto desconstruiu a visão estritamente
teórica que possuíamos sobre o desenvolvimento de software. O maior desafio não residiu na
codificação técnica, mas sim na adaptação da interface para garantir a acessibilidade de usuários com
pouca familiaridade digital, o que exigiu empatia e constante escuta ativa. Compreendemos o papel
social transformador da tecnologia quando aplicada à resolução de problemas reais de comunidades
historicamente marginalizadas. Essa interação quebrou os muros da universidade, humanizou nossa
formação técnica e evidenciou a urgência de cientistas e engenheiros comprometidos com o bem-estar
e o desenvolvimento social.

\vspace{0.7cm}
\noindent
\textbf{Referências}
[1]ASSOCIAÇÃO BRASILEIRA DE NORMAS TÉCNICAS. NBR 6022: Informação e documentação: artigo em
publicação periódica técnica e/ou científica: apresentação. Rio de Janeiro, 2018.

[2]SANTOS, M. A. A inclusão social dos catadores de materiais recicláveis através da tecnologia. 2. ed. São
Paulo: Editora Expressão Popular, 2022.

[3]UNIVERSIDADE FEDERAL DE SÃO PAULO. Projeto Pedagógico Institucional: a extensão universitária na
Unifesp. São Paulo: Unifesp, 2024. Disponível em: <https://www.unifesp.br>. Acesso em: 15 mar. 2026.

\vspace{0.7cm}
\chapter{Banco de dados}

\chapter{UC teste 2}

\section{Projeto - turma teste}
\noindent
\textbf{Autores}
Aluno Teste2

\vspace{0.7cm}
\noindent
\textbf{Disciplina e Curso}
BCT

\vspace{0.7cm}
\noindent
\textbf{Programa PEPICT vinculado}
Educação e Sustentabilidade

\vspace{0.7cm}
\noindent
\textbf{Objetivos}
objetivos teste

\vspace{0.7cm}
\noindent
\textbf{Metodologia}
metodologia teste

\vspace{0.7cm}
\noindent
\textbf{Resultados e Produtos}
resultado teste

\vspace{0.7cm}
\noindent
\textbf{Relação com os ODS}
ods teste

\vspace{0.7cm}
\noindent
\textbf{Reflexão Crítica}
teste

\vspace{0.7cm}
\noindent
\textbf{Referências}
teste

\vspace{0.7cm}


\end{document}